% GUID v6, Primer P4 (label box:cq): classical-quantum dynamics and collapse. Background only; it states no result of the paper.
% Place it before "Consistent classical-quantum dynamics" (marker %PRIMER:cq in part1.tex): % GUID v6, Primer P4 (label box:cq): classical-quantum dynamics and collapse. Background only; it states no result of the paper.
% Place it before "Consistent classical-quantum dynamics" (marker %PRIMER:cq in part1.tex): % GUID v6, Primer P4 (label box:cq): classical-quantum dynamics and collapse. Background only; it states no result of the paper.
% Place it before "Consistent classical-quantum dynamics" (marker %PRIMER:cq in part1.tex): % GUID v6, Primer P4 (label box:cq): classical-quantum dynamics and collapse. Background only; it states no result of the paper.
% Place it before "Consistent classical-quantum dynamics" (marker %PRIMER:cq in part1.tex): \input{primers/p4-cq}
\begin{primer}{Classical-quantum dynamics and collapse}{box:cq}
\setlength{\parskip}{3pt plus 1pt}%
\emph{The classical-quantum (CQ) master equation.} A consistent coupling keeps the joint description linear, positive and
normalised, as probability theory demands. The joint state is a family $\varrho(z)$ of positive operators, one for each value
$z$ of the classical variables (here the frame), normalised as $\int\dd z\,\Tr\varrho(z)=1$~\cite{oppenheim}. For one classical variable
$z$ driven by one quantum observable $L$ (for gravity, the mass density), continuous Markovian dynamics has the form~\cite{osswd}
\[
\partial_t\varrho=-\frac i\hbar[H,\varrho]+D_0\Big(L\varrho L-\tfrac12\{L^2,\varrho\}\Big)-\frac{D_1}2\,\partial_z\{L,\varrho\}
+\frac{D_2}2\,\partial_z^2\varrho ,
\]
with three coefficients beyond quantum mechanics:
\begin{itemize}[nosep,leftmargin=*]
\item $D_0$, \emph{decoherence}: the quantum system loses coherence between states that source the classical field differently (for
gravity, different mass distributions);
\item $D_1$, \emph{back-reaction}: a drift of $z$ driven by the quantum state, that is, how matter pushes the frame;
\item $D_2$, \emph{diffusion}: random kicks of $z$, that is, noise in the frame.
\end{itemize}
Complete positivity requires $D_1^2\le4D_0D_2$ here, and $D_0>0$ wherever $D_1\ne0$; Oppenheim et al.\ give the general matrix form of
this trade-off~\cite[Eq.~(23)]{osswd}, which is Eq.~\eqref{eq:tradeoff} of this paper. In words: \emph{a push implies noise}. If quantum matter pushes a
classical field at all, the matter decoheres and the field diffuses, and the product of the two is bounded below by the square of the push,
$D_0D_2\ge D_1^2/4$.

\emph{Measurement and feedback.} An older family of consistent models makes gravity a channel. Each mass's position, or its mass density,
is weakly and continuously measured, and a classical force computed from the record acts back on the other masses (Kafri, Taylor and
Milburn, KTM~\cite{ktm}; Tilloy and Di\'osi~\cite{tilloydiosi}). Newton's law then holds on average, the measurement decoheres
superpositions, and the feedback adds noise. The masses talk only through a classical record, so the channel is LOCC
(Primer~\ref{box:qi}) and cannot entangle them. A weaker measurement decoheres less but gives a noisier record and so more heating; KTM find
both effects at least of the order of the gravitational coupling. The least decoherence of two masses over all measurement strengths is
what this paper calls the \emph{KTM minimum}, $\omega/2$, where $\omega$ is the rate at which a coherent field would build up the
entangling phase.

\emph{Gravity-related collapse.} Di\'osi and Penrose proposed that a superposition of two mass distributions $\rho_1$, $\rho_2$ decays at
the rate $\lambda=E_\Delta/\hbar$, where
\[
E_\Delta=\frac G2\int\!\!\int\dd^3x\,\dd^3y\,\frac{[\rho_1(\mathbf x)-\rho_2(\mathbf x)][\rho_1(\mathbf y)-\rho_2(\mathbf y)]}
{|\mathbf x-\mathbf y|}
\]
is the gravitational self-energy of their difference~\cite{diosi87,penrose96,bassi2013}. For a body of mass $M$ and size $R$ displaced by
more than its size, $\lambda\approx GM^2/\hbar R$: the \emph{Di\'osi--Penrose (DP) rate}. For point-like nuclei $E_\Delta$ diverges, so the
mass density must be smeared over a length $R_0$. The same noise makes charged matter radiate. An underground germanium detector sees no
such radiation, which excludes smearing at nuclear size and requires $R_0$ of order $10^{-10}$\,m or more~\cite{donadi,figurato}. In the
Newtonian limit of CQ gravity the DP kernel is one admissible choice of the decoherence term~\cite{osswd}.
\end{primer}

\begin{primer}{Classical-quantum dynamics and collapse}{box:cq}
\setlength{\parskip}{3pt plus 1pt}%
\emph{The classical-quantum (CQ) master equation.} A consistent coupling keeps the joint description linear, positive and
normalised, as probability theory demands. The joint state is a family $\varrho(z)$ of positive operators, one for each value
$z$ of the classical variables (here the frame), normalised as $\int\dd z\,\Tr\varrho(z)=1$~\cite{oppenheim}. For one classical variable
$z$ driven by one quantum observable $L$ (for gravity, the mass density), continuous Markovian dynamics has the form~\cite{osswd}
\[
\partial_t\varrho=-\frac i\hbar[H,\varrho]+D_0\Big(L\varrho L-\tfrac12\{L^2,\varrho\}\Big)-\frac{D_1}2\,\partial_z\{L,\varrho\}
+\frac{D_2}2\,\partial_z^2\varrho ,
\]
with three coefficients beyond quantum mechanics:
\begin{itemize}[nosep,leftmargin=*]
\item $D_0$, \emph{decoherence}: the quantum system loses coherence between states that source the classical field differently (for
gravity, different mass distributions);
\item $D_1$, \emph{back-reaction}: a drift of $z$ driven by the quantum state, that is, how matter pushes the frame;
\item $D_2$, \emph{diffusion}: random kicks of $z$, that is, noise in the frame.
\end{itemize}
Complete positivity requires $D_1^2\le4D_0D_2$ here, and $D_0>0$ wherever $D_1\ne0$; Oppenheim et al.\ give the general matrix form of
this trade-off~\cite[Eq.~(23)]{osswd}, which is Eq.~\eqref{eq:tradeoff} of this paper. In words: \emph{a push implies noise}. If quantum matter pushes a
classical field at all, the matter decoheres and the field diffuses, and the product of the two is bounded below by the square of the push,
$D_0D_2\ge D_1^2/4$.

\emph{Measurement and feedback.} An older family of consistent models makes gravity a channel. Each mass's position, or its mass density,
is weakly and continuously measured, and a classical force computed from the record acts back on the other masses (Kafri, Taylor and
Milburn, KTM~\cite{ktm}; Tilloy and Di\'osi~\cite{tilloydiosi}). Newton's law then holds on average, the measurement decoheres
superpositions, and the feedback adds noise. The masses talk only through a classical record, so the channel is LOCC
(Primer~\ref{box:qi}) and cannot entangle them. A weaker measurement decoheres less but gives a noisier record and so more heating; KTM find
both effects at least of the order of the gravitational coupling. The least decoherence of two masses over all measurement strengths is
what this paper calls the \emph{KTM minimum}, $\omega/2$, where $\omega$ is the rate at which a coherent field would build up the
entangling phase.

\emph{Gravity-related collapse.} Di\'osi and Penrose proposed that a superposition of two mass distributions $\rho_1$, $\rho_2$ decays at
the rate $\lambda=E_\Delta/\hbar$, where
\[
E_\Delta=\frac G2\int\!\!\int\dd^3x\,\dd^3y\,\frac{[\rho_1(\mathbf x)-\rho_2(\mathbf x)][\rho_1(\mathbf y)-\rho_2(\mathbf y)]}
{|\mathbf x-\mathbf y|}
\]
is the gravitational self-energy of their difference~\cite{diosi87,penrose96,bassi2013}. For a body of mass $M$ and size $R$ displaced by
more than its size, $\lambda\approx GM^2/\hbar R$: the \emph{Di\'osi--Penrose (DP) rate}. For point-like nuclei $E_\Delta$ diverges, so the
mass density must be smeared over a length $R_0$. The same noise makes charged matter radiate. An underground germanium detector sees no
such radiation, which excludes smearing at nuclear size and requires $R_0$ of order $10^{-10}$\,m or more~\cite{donadi,figurato}. In the
Newtonian limit of CQ gravity the DP kernel is one admissible choice of the decoherence term~\cite{osswd}.
\end{primer}

\begin{primer}{Classical-quantum dynamics and collapse}{box:cq}
\setlength{\parskip}{3pt plus 1pt}%
\emph{The classical-quantum (CQ) master equation.} A consistent coupling keeps the joint description linear, positive and
normalised, as probability theory demands. The joint state is a family $\varrho(z)$ of positive operators, one for each value
$z$ of the classical variables (here the frame), normalised as $\int\dd z\,\Tr\varrho(z)=1$~\cite{oppenheim}. For one classical variable
$z$ driven by one quantum observable $L$ (for gravity, the mass density), continuous Markovian dynamics has the form~\cite{osswd}
\[
\partial_t\varrho=-\frac i\hbar[H,\varrho]+D_0\Big(L\varrho L-\tfrac12\{L^2,\varrho\}\Big)-\frac{D_1}2\,\partial_z\{L,\varrho\}
+\frac{D_2}2\,\partial_z^2\varrho ,
\]
with three coefficients beyond quantum mechanics:
\begin{itemize}[nosep,leftmargin=*]
\item $D_0$, \emph{decoherence}: the quantum system loses coherence between states that source the classical field differently (for
gravity, different mass distributions);
\item $D_1$, \emph{back-reaction}: a drift of $z$ driven by the quantum state, that is, how matter pushes the frame;
\item $D_2$, \emph{diffusion}: random kicks of $z$, that is, noise in the frame.
\end{itemize}
Complete positivity requires $D_1^2\le4D_0D_2$ here, and $D_0>0$ wherever $D_1\ne0$; Oppenheim et al.\ give the general matrix form of
this trade-off~\cite[Eq.~(23)]{osswd}, which is Eq.~\eqref{eq:tradeoff} of this paper. In words: \emph{a push implies noise}. If quantum matter pushes a
classical field at all, the matter decoheres and the field diffuses, and the product of the two is bounded below by the square of the push,
$D_0D_2\ge D_1^2/4$.

\emph{Measurement and feedback.} An older family of consistent models makes gravity a channel. Each mass's position, or its mass density,
is weakly and continuously measured, and a classical force computed from the record acts back on the other masses (Kafri, Taylor and
Milburn, KTM~\cite{ktm}; Tilloy and Di\'osi~\cite{tilloydiosi}). Newton's law then holds on average, the measurement decoheres
superpositions, and the feedback adds noise. The masses talk only through a classical record, so the channel is LOCC
(Primer~\ref{box:qi}) and cannot entangle them. A weaker measurement decoheres less but gives a noisier record and so more heating; KTM find
both effects at least of the order of the gravitational coupling. The least decoherence of two masses over all measurement strengths is
what this paper calls the \emph{KTM minimum}, $\omega/2$, where $\omega$ is the rate at which a coherent field would build up the
entangling phase.

\emph{Gravity-related collapse.} Di\'osi and Penrose proposed that a superposition of two mass distributions $\rho_1$, $\rho_2$ decays at
the rate $\lambda=E_\Delta/\hbar$, where
\[
E_\Delta=\frac G2\int\!\!\int\dd^3x\,\dd^3y\,\frac{[\rho_1(\mathbf x)-\rho_2(\mathbf x)][\rho_1(\mathbf y)-\rho_2(\mathbf y)]}
{|\mathbf x-\mathbf y|}
\]
is the gravitational self-energy of their difference~\cite{diosi87,penrose96,bassi2013}. For a body of mass $M$ and size $R$ displaced by
more than its size, $\lambda\approx GM^2/\hbar R$: the \emph{Di\'osi--Penrose (DP) rate}. For point-like nuclei $E_\Delta$ diverges, so the
mass density must be smeared over a length $R_0$. The same noise makes charged matter radiate. An underground germanium detector sees no
such radiation, which excludes smearing at nuclear size and requires $R_0$ of order $10^{-10}$\,m or more~\cite{donadi,figurato}. In the
Newtonian limit of CQ gravity the DP kernel is one admissible choice of the decoherence term~\cite{osswd}.
\end{primer}

\begin{primer}{Classical-quantum dynamics and collapse}{box:cq}
\setlength{\parskip}{3pt plus 1pt}%
\emph{The classical-quantum (CQ) master equation.} A consistent coupling keeps the joint description linear, positive and
normalised, as probability theory demands. The joint state is a family $\varrho(z)$ of positive operators, one for each value
$z$ of the classical variables (here the frame), normalised as $\int\dd z\,\Tr\varrho(z)=1$~\cite{oppenheim}. For one classical variable
$z$ driven by one quantum observable $L$ (for gravity, the mass density), continuous Markovian dynamics has the form~\cite{osswd}
\[
\partial_t\varrho=-\frac i\hbar[H,\varrho]+D_0\Big(L\varrho L-\tfrac12\{L^2,\varrho\}\Big)-\frac{D_1}2\,\partial_z\{L,\varrho\}
+\frac{D_2}2\,\partial_z^2\varrho ,
\]
with three coefficients beyond quantum mechanics:
\begin{itemize}[nosep,leftmargin=*]
\item $D_0$, \emph{decoherence}: the quantum system loses coherence between states that source the classical field differently (for
gravity, different mass distributions);
\item $D_1$, \emph{back-reaction}: a drift of $z$ driven by the quantum state, that is, how matter pushes the frame;
\item $D_2$, \emph{diffusion}: random kicks of $z$, that is, noise in the frame.
\end{itemize}
Complete positivity requires $D_1^2\le4D_0D_2$ here, and $D_0>0$ wherever $D_1\ne0$; Oppenheim et al.\ give the general matrix form of
this trade-off~\cite[Eq.~(23)]{osswd}, which is Eq.~\eqref{eq:tradeoff} of this paper. In words: \emph{a push implies noise}. If quantum matter pushes a
classical field at all, the matter decoheres and the field diffuses, and the product of the two is bounded below by the square of the push,
$D_0D_2\ge D_1^2/4$.

\emph{Measurement and feedback.} An older family of consistent models makes gravity a channel. Each mass's position, or its mass density,
is weakly and continuously measured, and a classical force computed from the record acts back on the other masses (Kafri, Taylor and
Milburn, KTM~\cite{ktm}; Tilloy and Di\'osi~\cite{tilloydiosi}). Newton's law then holds on average, the measurement decoheres
superpositions, and the feedback adds noise. The masses talk only through a classical record, so the channel is LOCC
(Primer~\ref{box:qi}) and cannot entangle them. A weaker measurement decoheres less but gives a noisier record and so more heating; KTM find
both effects at least of the order of the gravitational coupling. The least decoherence of two masses over all measurement strengths is
what this paper calls the \emph{KTM minimum}, $\omega/2$, where $\omega$ is the rate at which a coherent field would build up the
entangling phase.

\emph{Gravity-related collapse.} Di\'osi and Penrose proposed that a superposition of two mass distributions $\rho_1$, $\rho_2$ decays at
the rate $\lambda=E_\Delta/\hbar$, where
\[
E_\Delta=\frac G2\int\!\!\int\dd^3x\,\dd^3y\,\frac{[\rho_1(\mathbf x)-\rho_2(\mathbf x)][\rho_1(\mathbf y)-\rho_2(\mathbf y)]}
{|\mathbf x-\mathbf y|}
\]
is the gravitational self-energy of their difference~\cite{diosi87,penrose96,bassi2013}. For a body of mass $M$ and size $R$ displaced by
more than its size, $\lambda\approx GM^2/\hbar R$: the \emph{Di\'osi--Penrose (DP) rate}. For point-like nuclei $E_\Delta$ diverges, so the
mass density must be smeared over a length $R_0$. The same noise makes charged matter radiate. An underground germanium detector sees no
such radiation, which excludes smearing at nuclear size and requires $R_0$ of order $10^{-10}$\,m or more~\cite{donadi,figurato}. In the
Newtonian limit of CQ gravity the DP kernel is one admissible choice of the decoherence term~\cite{osswd}.
\end{primer}
